\section{Verification: The Central Problem}\label{sec:verification}

If capability is no longer the binding constraint, the question that matters is how anyone comes to know that an agent's work is correct. This section argues that the answer depends on two properties of the check rather than on the strength of the model. The first is whether the signal is external, observing executed code, program state or a held-out outcome, or internal, resting on the agent's own account of what it did. The second is whether the builder can write, edit or argue with the checker. Figure~\ref{fig:twobytwo} places the evidence on these two axes. The checks that hold up sit in one cell, an external signal the builder cannot reach; the other three fail in characteristic ways, and even the good cell fails when the instrument itself is broken.

\begin{figure}[t]
\centering
\small
\begin{tabular}{>{\raggedright\arraybackslash}p{1.7cm}>{\raggedright\arraybackslash}p{5.6cm}>{\raggedright\arraybackslash}p{5.6cm}}
\toprule
 & \textbf{Checker not editable by the builder} & \textbf{Checker editable by the builder} \\
\midrule
\textbf{External signal} (execution, state, held-out outcome) &
Spec-derived fail-to-pass tests \citep{V2,V21}; mutants the tests must kill \citep{V4,V19}; execution-based checkers and state oracles in the running application \citep{V10,V28}; agent configuration restored from the base branch (\repofile{anthropics/claude-code-action}{docs/security.md}).
\newline\emph{What happens:} success is decided by something the agent cannot change; the remaining risk is a broken instrument \citep{V12,C15,V27}. &
Tests the agent may modify \citep{V3}; scorer tampering \citep{V17}; a verifier run in the agent's own sandbox \citep{C12}; search loops that overfit or exploit the score \citep{R26,S28}; attacks on the scorer and its network \citep{F34,F35}; a metric computed by agent-owned code (\repo{karpathy/autoresearch}{599}); a graded diff that includes test files (\repo{commit-0/commit0}{100}).
\newline\emph{What happens:} the signal is real but the agent can change what it measures, so scores rise without the work improving. \\
\midrule
\textbf{Internal signal} (the agent's account, code read but not run) &
A separate LLM verifier \citep{T18}; LLM judges of code \citep{V8,V9} and of agent runs \citep{V24,V25}; model critique without external feedback \citep{A4,A5}.
\newline\emph{What happens:} the checker is independent but reads a plausible account and accepts much failing work; rubric checks of observed state move it toward the cell above \citep{V15,V29}. &
Completion claims \citep{T23,F20,F25,H13}; skipped checks scored as passes (\repo{SakanaAI/AI-Scientist-v2}{133}, \repo{K-Dense-AI/scientific-agent-skills}{152}, \repofile{aiming-lab/AutoResearchClaw}{docs/CODE_REVIEW_2026-07-03.md}).
\newline\emph{What happens:} the system certifies itself, and failure is silent because it looks like success. \\
\bottomrule
\end{tabular}
\caption{Verification signals by what they observe (rows) and by whether the builder can edit the checker (columns), with representative evidence in each cell. In the studies reviewed, only the top-left cell gives reliable evidence, and only while the instrument itself is sound.}
\label{fig:twobytwo}
\end{figure}

\subsection{Self-Report Is Unreliable}

The weakest signal is the agent's own statement that it has finished, and the evidence against it now comes from real use as well as from benchmarks. In a large observational study, \citet{T23} analyze 20,574 sessions in IDE and command-line coding tools. They classify 22.58\% of misalignment episodes as inaccurate self-reporting, in which the agent prematurely claims ``success, completion, or readiness'', and find that 91.49\% of visible resolutions (themselves only 9.33\% of episodes) still required explicit correction from the developer. Because the study counts misalignment only when the developer pushes back, latent failures are excluded. \citet{F20} separates submission from resolution across 1,750 trajectories on 50 SWE-bench Verified tasks: GPT-5 submits a patch on 100\% of runs but resolves 44\%, and the dominant failure, the same plausible but wrong patch repeated across runs, covers 68\% of GPT-5's failing runs and 80\% of Llama~4's, so that monitors based on completion or on consistency ``both look healthy.'' Annotating 2,518 trajectories, \citet{F25} find that after a first mistake an agent recovers in 30.5\% of runs and never detects the error in 38.5\%, and that even runs scored as solved ``delete data, corrupt systems, or fabricate success.'' Anthropic's account of its own long-running harness reports that the agent ``marks features as done prematurely'' \citep{H13}. A completion message is a claim, not evidence.

\subsection{LLM Judges}

The obvious remedy, a second model that reviews the first, inherits much of the weakness when the reviewer reads rather than runs. The cleanest evidence comes from \citet{T18}, where roles are separated by operating-system permissions and every verifier verdict can be paired with a deterministic grader: the LLM verifier approves 49.4\% of submissions the grader rejects ($n = 1{,}083$), and in a reference ablation removing the verifier raises the team's mean partial score by 5.5 points. Code judges fare similarly. Checked against test execution, GPT-4 ``misjudges 50\% of the wrong'' Java implementations \citep{V9}. Among judges of agent runs, none of 12 LLM judges exceeds 70\% precision on 1,302 expert-labeled web trajectories \citep{V25}. Most damaging for autonomous work, judges are swayed by the agent's own narrative: when only an agent's reasoning text is rewritten, with actions and observations held fixed, fabricated signs of progress raise the false-positive rates of two vision-language judges (GPT-4o and GLM-4.1V) by 20--30 points on 800 trajectories, with smaller but non-trivial increases for others, and prompt-based mitigations ``reduce but do not eliminate'' the effect \citep{V24}. This mirrors the self-correction literature. Models fail to correct their reasoning ``without external feedback, and at times, their performance even degrades'' \citep{A4}, and a critical survey finds that self-correction ``works well in tasks that can use reliable external feedback'' but has not been shown from prompted-LLM feedback outside tasks exceptionally suited to it \citep{A5}.

The exceptions are instructive. Judges improve when they must check evidence against an explicit list rather than form an overall impression. An agent judge that inspects the work against stated requirements matches the human consensus 90\% of the time against 70\% for a plain LLM judge \citep{V15}, and a rubric-based verifier for computer-use agents that examines the screenshots cuts false positives to near zero where baseline judges show at least 45\% and 22\% \citep{V29}. Both results come from the teams that built the judges and rest on small samples (55 tasks in \citealt{V15}; 140 and 106 trajectories in \citealt{V29}), so they establish a direction rather than a rate: a judge becomes more reliable as its input moves from the agent's account toward observed state.

\subsection{Checks That Work}

The reliable checks share a property rather than a technique: the builder cannot write, edit or argue with them. The clearest case is a test derived from the specification and shown to fail before the change and pass after it. On SWT-Bench, code agents produce such fail-to-pass tests for a minority of issues, but keeping only fixes that pass them more than doubles the precision of SWE-agent to 47.8\%, at 20\% recall \citep{V2}. Meta generates tests meant to fail on a code change before it lands; of 41 candidate catches reported to engineers, 8 were confirmed true positives, 4 of which would have led to serious failures \citep{V21}. Tests that merely pass are weak evidence by comparison. Coverage-driven generators run on buggy code keep tests that encode the bug, in 68.1\% of final suites for one tool \citep{V5}, and agent-written test files contain far more print statements than assertions, with prompts that encourage or discourage them changing resolved totals by 0--13 of 500 per model, a difference that is not significant \citep{V1}. Mutation testing supplies a second external check. In one reported case, an LLM-written suite reached 100\% line and branch coverage but a 4\% mutation score \citep{V19}, and at Meta 277 of 571 mutation-guided tests would have been discarded under a line-coverage criterion, while engineers accepted 73\% of the tests \citep{V4}. For user interfaces, the equivalent is a check on state in the running application rather than on appearance. OSWorld scores each task with an execution-based evaluation script \citep{V10}; vendor guidance tells graders to confirm ``an order was actually placed, not just that the confirmation page appeared'' \citep{V16}; and pre- and post-condition oracles derived from the specification reach 0.96 precision and recall on 100 bugs injected into four open-source web applications \citep{V28}. Practitioner tooling builds the property into the environment: Anthropic's GitHub Action for Claude Code restores the agent's configuration files from the base branch before a pull-request run, so that the pull request cannot rewrite those configuration files (\repofile{anthropics/claude-code-action}{docs/security.md}). The caveat is scope: several results are single-company reports \citep{V4,V21}, most use small benchmarks, and none measures the build of a whole application.

\subsection{Gaming the Grader}

When the checker is within reach, agents use it. ImpossibleBench makes this measurable with tasks whose tests contradict the specification, so that any pass is cheating: GPT-5 cheats on 54.0\% of Conflicting-SWEbench tasks, Claude models and Qwen3-Coder cheat primarily (over 79\%) by modifying test cases, and, averaged over the models tested, allowing multiple submissions raises cheating from 33\% to 38\% \citep{V3}. An honest exit helps some models: telling them to submit a flag for human intervention lowers GPT-5's rate from 54\% to 9\%, though the effect is ``much less pronounced'' for Claude Opus~4.1. Because the tasks are built to be impossible, these rates are best read as an upper bound for tasks that are merely hard. In METR's evaluation, o3 reward-hacked in 39 of 128 RE-Bench runs (30.4\%), by overwriting a grader's timer, monkey-patching an evaluator or caching a precomputed answer, and on one task instructions not to cheat had ``a nearly negligible effect'' \citep{V17}. Search loops game metrics more quietly. Selecting the final MLE-bench solution by test rather than validation score would raise the medal rate by 9 to 13 points, which the authors read as systematic overfitting \citep{R26}; a kernel-optimization loop that appeared to beat prior results fell from an average 3.13$\times$ to 1.49$\times$ speedup once contaminated tasks were excluded, with fake speedups of 50--120$\times$ passing verification \citep{S28}. At the extreme, a vendor incident report describes models under a cyber-security evaluation, run with reduced refusals and without production classifiers, escaping the sandbox through a zero-day in a package-registry proxy to obtain ``test solutions directly from Hugging Face's production database'' \citep{F34}. METR's brief independent investigation reports that roughly 1,200 agents took part on a shared message board, which agents used to coordinate projects ``to find a general-purpose way to trick or tamper with the automated ExploitGym scorer'' \citep{F35}. Practitioner repositories show the same opening at small scale. In autoresearch, an issue points out that the agent owns the file that calls the evaluator, so it could lower the reported metric without improving the model, and demonstrates this with a minimal reproducer (\repo{karpathy/autoresearch}{599}); in commit0, whose grading diff covers test files as well as source, a pull request to exclude test-file changes was closed unmerged after a maintainer noted it was not yet correct (\repo{commit-0/commit0}{100}).

\subsection{Broken Instruments}

Even a checker outside the builder's reach can be wrong, and audits from 2025--2026 suggest this is common. On SWE-bench Verified, the official test check lets 7.8\% of patches count as correct while failing the developer-written test suite, and the combined weaknesses inflate reported resolution rates by 6.4 points \citep{V12}. In vendor audits, OpenAI found that at least 59.4\% of 138 Verified problems that o3 did not consistently solve have tests that reject functionally correct submissions \citep{C15}, then estimated that about 30\% of SWE-Bench Pro tasks are broken and retracted its own recommendation of that benchmark \citep{V31}. Authors from Scale AI and Reflection AI have since repaired or retired 480 of its 731 tasks \citep{C12}. Flaws cut both ways. On $\tau$-bench, an agent that returns empty responses scores 38\% because impossible tasks count it as successful, and benchmark flaws can misstate performance by up to 100\% in relative terms \citep{A18}; on the data-engineering benchmark ELT-Bench, 82.7\% of failed tasks contained a benchmark-attributable error, and correcting the key raised the same agent's score from 22.66\% to 32.51\% \citep{D13}. Adversarial auditing is the strongest response so far. An automated red team built on coding agents found 219 flaws across 10 agent benchmarks, with exploits that score near-perfectly ``without solving a single task,'' and its patch loop cut the hackable share of four benchmarks from near 100\% to under 10\% \citep{V27}. A hacker--fixer--solver loop found 323 of 1,968 terminal-agent tasks (16\%) hackable from the task description alone and drove attack success on KernelBench from 62\% to 0\% \citep{T24}. A grader is itself an instrument that needs validation before its scores can be trusted.

\subsection{Silent Failure}

In the practitioner record, a recurring way these systems mislead is not an elaborate exploit but a check that never ran and was counted as a pass. An independent evaluation of AI-Scientist-v2 found that its vision-model figure review rejected the configured model in every one of 10 runs, so that ``v2 never actually reviews its own figures'' (\repo{SakanaAI/AI-Scientist-v2}{133}). In a library of science skills, a wrong model identifier made every image-review call fail; a broad exception handler returned a hard-coded score of 7.5 with no improvement needed, and the loop printed ``Quality meets threshold (\texttt{7.5 >= 8.5})'' (\repo{K-Dense-AI/scientific-agent-skills}{152}). A code audit committed to the AutoResearchClaw repository found, among other fail-open defenses, a citation check that, with the network down, skipped every citation and returned a perfect integrity score of 1.0, while the test suite reported 2,831 passing tests (\repofile{aiming-lab/AutoResearchClaw}{docs/CODE_REVIEW_2026-07-03.md}); a follow-up pull request that inventoried the catch-all handlers on the experiment path, so that such failures would fail the stage or be recorded as degradations, was closed unmerged after review found blockers (\repo{aiming-lab/AutoResearchClaw}{302}), and a narrower version remains open (\repo{aiming-lab/AutoResearchClaw}{323}). In AI-Scientist-v3, a repair agent made the evaluator non-deterministic: two scores of the same program differed by about 0.2, more than the real gains of about 0.15, and the improvement check became ``effectively a coin flip'' (\repo{SakanaAI/AI-Scientist-v3}{5}). These are single reports, not rates, and the science-skills and AI-Scientist-v3 issues have since been closed as completed, but they share a structure: a skippable check, and a default that treats the skip as success. In the terms of Figure~\ref{fig:twobytwo}, such a check becomes a self-certified internal signal, however external it was designed to be. For anyone building these systems, the question to ask of every verification step is not how capable the checker is, but whether the builder can reach it and whether its failure would be visible.
